% TOP_PROBLEM: {"id": 3263, "title": "Caccetta-Häggkvist Conjecture", "category_id": 3, "category": {"id": 3, "name": "graph_theory", "display_name": "Graph Theory", "description": "Problems involving graphs, networks, and their properties.", "slug": "graph-theory", "order_index": 3, "created_at": "2026-07-31T15:26:25.671Z"}, "status": "open", "problem_number": "OPG-46385", "set_id": 12}
% FIRST_POSTED: 2026-10-02
% ATTEMPT_STATUS: unresolved
% UnsolvedMath ID 3263; source catalogue code OPG-46385
% !TeX program = lualatex
% GPT-6 Astra Ultra research attempt; independently checked partial work, 2026-10-02.
\documentclass[11pt,a4paper]{article}
\usepackage[margin=25mm]{geometry}
\usepackage{fontspec}
\setmainfont{lmroman10-regular.otf}[BoldFont=lmroman10-bold.otf,
 ItalicFont=lmroman10-italic.otf,BoldItalicFont=lmroman10-bolditalic.otf]
\usepackage{amsmath,amssymb,mathtools}
\usepackage{unicode-math}
\setmathfont{latinmodern-math.otf}
\AtBeginDocument{\let\setminus\smallsetminus}
\usepackage{xurl,hyperref}
\hypersetup{colorlinks=true,urlcolor=blue}
\setcounter{secnumdepth}{0}
\setlength{\parindent}{0pt}
\setlength{\parskip}{5pt}
\setlength{\emergencystretch}{3em}
\begin{document}
\section*{Caccetta-Häggkvist Conjecture}\label{3263}
{\small UnsolvedMath ID 3263 · OPG-46385 · Graph Theory · OpenGarden}
\subsection{Definitions and mathematical statement}
Let $n\ge2$ and $1\le r\le n-1$ be integers. A simple digraph on $n$
vertices has no loops and has at most one arc for each ordered pair of
distinct vertices; opposite arcs are permitted. Its minimum outdegree
$\delta^+(D)$ is the smallest number of distinct outneighbours of a
vertex. A directed cycle of length $\ell\ge2$ has $\ell$ distinct
vertices and follows the arc directions around them; opposite arcs
therefore constitute a cycle of length two.

The Caccetta--Häggkvist conjecture is
\[
 \delta^+(D)\ge r\quad\Longrightarrow\quad
 D\text{ has a directed cycle of length at most }
                         \left\lceil\frac nr\right\rceil.
\]
The ceiling is the least integer at least $n/r$. The quantifiers cover
every finite simple digraph satisfying the degree hypothesis. No
indegree bound, regularity, or connectedness assumption is imposed.
Because every vertex has an outgoing arc, following arcs eventually
produces a cycle; the problem is the specific length guarantee.

For example, when $r\ge n/3$ the assertion asks for a cycle of length
two or three. When the digraph is an orientation and hence has no
2-cycle, this becomes a directed-triangle assertion. The parameter $n$
counts vertices of the whole digraph and $r$ is a uniform lower bound,
not an average degree. Replacing a cycle by a closed walk would weaken
the formulation, although any shortest directed closed walk contains
a simple directed cycle.
\subsection{Short English statement}
Does an n-vertex digraph of minimum outdegree r always contain a directed cycle of length at most the ceiling of n/r?
\subsection{Research attempt: sharp circulants and limits of regularization}
\paragraph{Scope and literature check.}
GPT-6 Astra Ultra, 2 October 2026. Opposite arcs are allowed and count
as a two-cycle. The hypothesis concerns minimum outdegree, not equal
in- and outdegrees. Ho\`ang and Reed [S3] proved the conjecture for
outdegree parameters four and five; Sullivan [S4] surveys the
conjecture and several related problems. Those references are
context, not a claim that their bounds represent the latest complete
state of this large literature. The work here proves boundary cases,
sharp examples for every parameter, and an explicit obstruction to
a tempting reduction to regular digraphs.

\paragraph{1. Two directly provable regimes.}
For $r=1$, repeatedly follow an outgoing arc. A repeated vertex
produces a simple directed cycle of length at most $n$, exactly the
required bound. For $r\ge n/2$, there are at least $nr$ arcs. If
there were no opposite pair, there would be at most $n(n-1)/2<nr$
arcs. Thus a two-cycle exists. Since $r\le n-1$, the value
$\lceil n/r\rceil$ in this regime is two, so this proves the target.
These arguments include the smallest allowed orders and do not
assume connectedness.

\paragraph{2. A valid reduction, followed by an invalid stronger one.}
The condensation of a digraph into its strongly connected components
is acyclic and has a sink component. All outgoing arcs of a vertex
in that component stay inside it, so its induced digraph retains
minimum outdegree at least $r$. If its order is $m\le n$, a cycle
of length at most $\lceil m/r\rceil$ would suffice for the original
problem. Thus it is enough to prove the conjecture for strongly
connected digraphs.

Deleting excess outgoing arcs also reduces to exactly $r$ arcs
leaving every vertex, without changing the vertex set. Such deletion
need not preserve strong connectivity; applying the sink-component
reduction again is legitimate. Neither operation guarantees equal
indegrees. For example, on $\mathbb Z/5\mathbb Z$ start with arcs
$x\to x+1,x+2$ and replace $0\to2$ by $0\to3$. The unit-step
Hamilton cycle remains, so this digraph is strong and every outdegree
is two, but its indegrees are $(2,2,1,3,2)$. It cannot have a spanning
subdigraph with every indegree and outdegree two: retaining two
outgoing arcs at each vertex forces retaining every arc. Therefore
regularizing by arc deletion is not a valid general reduction.

\paragraph{3. Exact sharpness for every $(n,r)$.}
Let $D_{n,r}$ have vertices $\mathbb Z/n\mathbb Z$ and arcs
$x\to x+s$ for $1\le s\le r$. Its outdegree is exactly $r$.
On a directed cycle of length $\ell$, select the step representative
in $\{1,\ldots,r\}$ for each arc. Their sum is a positive multiple
of $n$, so $\ell r\ge n$. Hence the directed girth is at least
$\lceil n/r\rceil$.

Put $h=\lceil n/r\rceil$. The steps
\[
 \underbrace{r,\ldots,r}_{h-1\text{ times}},
                    n-(h-1)r
\]
all lie between one and $r$ and sum to $n$. Their proper partial
sums lie strictly between zero and $n$, and are strictly increasing.
They therefore define a simple directed $h$-cycle. The girth is
exactly $h$. This verifies the ceiling and shows that any smaller
universal target would be false, including when two-cycles occur.

\paragraph{4. Why counting an out-tree is insufficient.}
High directed girth does not make different outgoing branches
disjoint. In $D_{10,2}$, which has directed girth five, vertices zero
and one both have vertex two as an outneighbour. A breadth-first
exploration can merge without closing any short directed cycle.
Thus multiplying the minimum outdegree at successive levels is
not a justified vertex lower bound; a directed analogue of an
undirected tree expansion cannot simply be assumed.

\paragraph{Verification, gap and final status.}
The script [S9] computes shortest directed cycles for every
$D_{n,r}$ with $2\le n\le30$ and $1\le r<n$, and verifies the
strong, unbalanced example. The first remaining gap is controlling
short cycles for arbitrary outregular digraphs when branches merge
and indegrees vary. No regularity theorem or expansion estimate
covering that gap is established. Full target: UNRESOLVED.
\subsection{Sources}
\begin{itemize}
\item[S1] ULAM.AI and the UnsolvedMath Contributors, supplied problems.json, record 3263 (OPG-46385), OpenGarden collection. Catalogue identity and source transcription; CC BY 4.0.
\url{https://huggingface.co/datasets/ulamai/UnsolvedMath}.
\item[S2] Open Problem Garden, Caccetta-Häggkvist Conjecture. Original problem and discussion; the mathematical formulation below is expanded from the supplied source record.
\url{http://www.openproblemgarden.org/op/caccetta_haggkvist_conjecture}.
\item[S3] C. T. Ho\`ang and B. Reed, \emph{A note on short cycles in diagraphs}, Discrete Mathematics 66 (1987), 103--107. \url{https://doi.org/10.1016/0012-365X(87)90122-1}.
\item[S4] Blair D. Sullivan, \emph{A Summary of Problems and Results related to the Caccetta--H\"aggkvist Conjecture} (2006). \url{https://arxiv.org/abs/math/0605646}.
\item[S9] GPT-6 Astra Ultra, exact finite checks accompanying this attempt (2 October 2026). Repository script, Python standard library only. \texttt{scripts/check\_attempt\_batch03\_digraphs\_2026\_10\_02.py}.
\end{itemize}
\end{document}
